\documentclass[12pt]{article}
\usepackage{ficha}

\begin{document}

\ficha{Proporcionalidade Direta e Escalas}

\instrucoes{Resolve os exercícios, apresentando todos os cálculos e indicando as unidades nas respostas.}

\secao{Reconhecer Proporcionalidade Direta}{Indica se as grandezas são diretamente proporcionais. Em caso afirmativo, indica a constante de proporcionalidade.}
\begin{exercicios}
  \item \begin{tabular}[c]{|c|c|c|c|c|}\hline
    Quantidade de pães & 1 & 2 & 5 & 8 \\\hline
    Preço (€) & 0,30 & 0,60 & 1,50 & 2,40 \\\hline
  \end{tabular}
  \item \begin{tabular}[c]{|c|c|c|c|c|}\hline
    Idade (anos) & 2 & 4 & 6 & 8 \\\hline
    Altura (cm) & 85 & 100 & 115 & 128 \\\hline
  \end{tabular}
  \item \begin{tabular}[c]{|c|c|c|c|c|}\hline
    Tempo (h) & 1 & 3 & 4 & 6 \\\hline
    Distância (km) & 80 & 240 & 320 & 480 \\\hline
  \end{tabular}
  \item O lado de um quadrado e o seu perímetro.
  \item O lado de um quadrado e a sua área.
  \item O número de bilhetes de cinema comprados e o preço total (todos os bilhetes custam o mesmo).
  \item A idade de uma pessoa e o seu peso.
  \item O número de litros de gasolina e o preço a pagar.
\end{exercicios}

\secao{Completar Tabelas}{As grandezas são diretamente proporcionais. Completa as tabelas.}
\begin{exercicios}
  \item \begin{tabular}[c]{|c|c|c|c|c|}\hline
    $x$ & 2 & 5 & \phantom{00} & 12 \\\hline
    $y$ & 6 & \phantom{00} & 24 & \phantom{00} \\\hline
  \end{tabular}
  \item \begin{tabular}[c]{|c|c|c|c|c|}\hline
    Massa (kg) & 0,5 & 2 & \phantom{00} & 7,5 \\\hline
    Preço (€) & 1,20 & \phantom{000} & 7,20 & \phantom{000} \\\hline
  \end{tabular}
  \item \begin{tabular}[c]{|c|c|c|c|c|}\hline
    Tempo (min) & 4 & 10 & \phantom{00} & 30 \\\hline
    Água (L) & 18 & \phantom{00} & 72 & \phantom{00} \\\hline
  \end{tabular}
  \item Numa proporcionalidade direta, $y = 4$ quando $x = 10$. Escreve $y$ em função de $x$ e calcula $y$ para $x = 25$.
  \item O gráfico de uma proporcionalidade direta é uma reta. Por que ponto passa sempre essa reta?
\end{exercicios}

\secao{Problemas de Proporcionalidade Direta}{Resolve.}
\begin{exercicios}
  \item 6 lápis custam 4,20 euros. Quanto custam 15 lápis?
  \item Um carro percorre 210\,km com 14 litros de gasolina. Quantos quilómetros percorre com 30 litros?
  \item Uma impressora imprime 45 páginas em 3 minutos. Quanto tempo demora a imprimir 120 páginas?
  \item Para fazer 12 bolos usam-se 3\,kg de farinha. Quantos bolos se fazem com 5\,kg?
  \item Um trabalhador ganha 68 euros por 8 horas de trabalho. Quanto ganha em 35 horas?
  \item Uma torneira deita 12 litros em 5 minutos. Quantos litros deita numa hora?
  \item Numa receita, para 4 pessoas usam-se 300\,g de arroz. Quanto arroz é preciso para 10 pessoas?
\end{exercicios}

\secao{Escalas}{Resolve. A escala é a razão entre a medida no desenho e a medida real.}
\begin{exercicios}
  \item Num mapa à escala $1:50\,000$, 1\,cm representa quantos metros reais? E quantos quilómetros?
  \item Num mapa à escala $1:200\,000$, duas cidades estão a 7\,cm. Qual é a distância real em km?
  \item Numa planta à escala $1:100$, uma sala mede 4,5\,cm por 3,2\,cm. Quais são as dimensões reais?
  \item Uma estrada de 18\,km está representada num mapa por 6\,cm. Qual é a escala do mapa?
  \item Um campo de futebol com 105\,m de comprimento foi desenhado com 21\,cm. Qual é a escala?
  \item Numa maqueta à escala $1:25$, um carro mede 16\,cm. Qual é o comprimento real?
  \item Um inseto com 4\,mm foi desenhado com 6\,cm. Qual é a escala do desenho?
  \item Que medida terá, num mapa à escala $1:25\,000$, uma distância real de 3,5\,km?
\end{exercicios}

\secao{Desafios}{Resolve os problemas.}
\begin{exercicios}
  \item Num mapa, 2\,cm representam 15\,km. Qual é a escala? A que distância real correspondem 9\,cm?
  \item Uma planta de uma casa está à escala $1:50$. Um quarto quadrado tem 16\,m$^2$ de área real. Que área tem na planta, em cm$^2$?
  \item Numa mistura de tinta usam-se 3 partes de azul para 5 de branco. Quantos litros de azul são necessários para 24 litros de mistura?
  \item O preço do combustível é diretamente proporcional à quantidade. 40 litros custam 70 euros. Escreve uma expressão para o custo $C$ de $x$ litros e calcula o custo de 55 litros.
  \item Duas grandezas $x$ e $y$ são diretamente proporcionais e $x + y = 36$. Se a constante de proporcionalidade $\dfrac{y}{x}$ é 3, determina $x$ e $y$.
  \item Num mapa à escala $1:300\,000$, um lago tem forma retangular com 2\,cm por 1,5\,cm. Qual é a área real do lago, em km$^2$?
\end{exercicios}

\end{document}
